\chapter{创新性、适用价值与发展方向}

\section{创新性分析}

VeriSpecOSLab 的创新集中在一套可以直接用于操作系统课程的平台。平台把学生设计、AI 修改、实际构建、QEMU 或开发板运行、教师复核和最终提交连接起来。每个环节都有明确输入和输出，学生能够解释自己的设计和代码，教师能够从报告回到原始规格与运行记录。

五类规格把系统方案、模块职责、跨模块接口、功能目标和跨模块修改分别记录。ModuleSpec 中的 \code{owns} 用于确定当前任务的修改范围，properties 与 checks 用于确定验收内容，SpecPatch 用于说明跨模块变化。规格由此同时参与 AI 实现和后续验证。

学生可以先通过问答助手理解课程概念和讨论设计，再亲自编写规格。平台检查规格结构、引用关系和验收映射，评审助手只读取文件并给出修改意见。学生完成修改并提交规格后，实现助手在独立临时工作目录中生成代码和配套测试。平台依据 Git 实际变化检查文件范围，再执行构建、公开测试、契约测试、固定种子 fuzz 和有限 trace 测试。通过检查的实现以单独提交进入学生仓库，并用 Run-ID 和 Spec-Hash 对应本次任务。

平台还统一描述宿主构建、QEMU 和开发板任务。每次运行记录代码版本、程序参数、环境变量、日志、产物和判定结果。开发板实验增加教学团队复核，自动运行记录与人工判断一起归档。xv6-spec 的 Lab 1--10 课程历史进一步说明同一机制可以覆盖从最小内核到 RISC-V 多核开发板的连续实验。

\section{与相关课程工具的关系}

\begin{table}[htbp]
  \centering
  \caption{相关工具与 VeriSpecOSLab 的结合方式}
  \begin{tabularx}{\textwidth}{p{3cm}p{4.1cm}X}
    \toprule
    工具或范式 & 主要能力 & VeriSpecOSLab 的结合方式 \\
    \midrule
    MIT xv6 教学 & 小型、可读的多核 RISC-V Unix-like 内核与实验 & 作为完整源码课程案例，并补充规格、AI 协作、运行证据和硬件实验流程 \\
    rCore 类教学 & 以 Rust 和分阶段实现组织操作系统学习 & 平台的规格与运行接口可以承载不同语言和内核结构的课程案例 \\
    通用编码 Agent & 阅读仓库、编辑文件和执行命令 & 增加课程问答、学生手写规格、修改范围、实际构建测试和可复查结果 \\
    SYSSPEC 与 SPECFS & 多部分规格、规格补丁和生成式系统方法 & 吸收规格分解与补丁思想，并用于操作系统课程和学生主导的实验流程 \\
    在线 OJ 与 Lab 平台 & 统一测试、课程管理和结果汇总 & 保留的 Portal 组件复用本地规格、Runner 和证据类型，当前学生流程由 CLI 完成 \\
    \bottomrule
  \end{tabularx}
\end{table}

\section{课程应用价值}

对学生，平台要求其在编码前说明设计，亲自修改规格，并在提交前解释代码、测试与运行结果。AI 提供概念问答、规格评审、代码实现和调试帮助，学生通过手写规格、Git 提交和实验复盘保留自己的学习过程。

对教师和助教，平台提供统一的规格模板、课程知识库、运行步骤、日志和提交材料。教师可以检查某次修改对应哪个 Lab、涉及哪些模块、执行了哪些测试，也可以在答辩中根据规格和代码提出更具体的问题。

对课程建设者，结构化配置、稳定 Spec ID、累计课程标签和可复现知识来源使课程内容能够共同版本化。xv6-spec 展示了一套从 Lab 1 到 Lab 10 的完整案例，新的课程可以复用平台流程并替换内核、工具链和验收方法。

\section{发展方向}

\subsection{扩大教学应用与学习评价}

暑期两节试讲发现了平台使用、前置知识和大型项目开发经验三方面的实际问题，并推动课程材料和学生流程完成调整。后续教学可以覆盖更多学生和完整 Lab 周期，记录环境准备时间、概念理解、代码解释能力、调试时间、测试覆盖和教师复查效率。课程评价将结合过程材料、代码结果与现场答辩，观察学生如何使用 AI 理解进程、内存、文件系统和设备驱动。

\subsection{更多课程案例}

后续可以增加 rCore、网络协议栈、数据库内核和编译器运行时等案例。每个案例提供自己的 DesignSpec、ModuleSpec、InterfaceSpec、构建目标和验收方法，共享 Agent、Runner、知识库、报告和提交能力。

\subsection{更丰富的验证方法}

现有 checks 可以继续接入 property-based testing、差分测试、模糊测试、有限状态模型检查和关键不变量验证。教师根据 Lab 难度选择相应方法，并在规格中写明验收目标。

\subsection{课程协作与教学资料}

课程教学团队可以共同维护规格模板、知识来源、运行步骤、测试目标和答辩材料，并通过稳定 Spec ID 与课程标签同步版本。保留的 Portal API、Web UI 和 worker 继续参与类型检查、构建和单元测试，为后续课程协作研究保存一致的接口基础。

\begin{keybox}[发展原则]
平台的扩展围绕课程实际使用展开。每项新能力都对应清晰的学生任务、教师验收方式和可复查运行记录。
\end{keybox}
